\begin{table}[htbp]
\centering
\caption{UK Country Composition: Progressively Restricting Exposure}
\label{tab:distance}
\begin{tabular}{lcccc}
\toprule
UK Composition & Coefficient & SE & N UK Regions & Observations \\
\midrule
All UK & 0.0250** & (0.0114) & 183 & 3,523 \\
Exclude NI & 0.0251** & (0.0113) & 172 & 3,523 \\
England only & 0.0254** & (0.0109) & 118 & 3,523 \\
Non-England UK & 0.0211* & (0.0118) & 65 & 3,523 \\
\bottomrule
\end{tabular}
\begin{minipage}{0.9\textwidth}
\footnotesize \textit{Notes:} Each row restricts the set of UK GADM2 regions used to construct the SCI exposure measure. All specifications include d\'epartement and quarter-year fixed effects with d\'epartement-clustered standard errors (93 clusters). England's 118 regions dominate UK SCI due to population size. If effects are driven only by England-specific confounders, restricting to non-England regions should attenuate coefficients. $^{***}p<0.01$, $^{**}p<0.05$, $^{*}p<0.1$.
\end{minipage}
\end{table}
